% Options for packages loaded elsewhere
\PassOptionsToPackage{unicode}{hyperref}
\PassOptionsToPackage{hyphens}{url}
%
\documentclass[
]{article}
\usepackage{amsmath,amssymb}
\usepackage{iftex}
\ifPDFTeX
  \usepackage[T1]{fontenc}
  \usepackage[utf8]{inputenc}
  \usepackage{textcomp} % provide euro and other symbols
\else % if luatex or xetex
  \usepackage{unicode-math} % this also loads fontspec
  \defaultfontfeatures{Scale=MatchLowercase}
  \defaultfontfeatures[\rmfamily]{Ligatures=TeX,Scale=1}
\fi
\usepackage{lmodern}
\ifPDFTeX\else
  % xetex/luatex font selection
\fi
% Use upquote if available, for straight quotes in verbatim environments
\IfFileExists{upquote.sty}{\usepackage{upquote}}{}
\IfFileExists{microtype.sty}{% use microtype if available
  \usepackage[]{microtype}
  \UseMicrotypeSet[protrusion]{basicmath} % disable protrusion for tt fonts
}{}
\makeatletter
\@ifundefined{KOMAClassName}{% if non-KOMA class
  \IfFileExists{parskip.sty}{%
    \usepackage{parskip}
  }{% else
    \setlength{\parindent}{0pt}
    \setlength{\parskip}{6pt plus 2pt minus 1pt}}
}{% if KOMA class
  \KOMAoptions{parskip=half}}
\makeatother
\usepackage{xcolor}
\usepackage[margin=1in]{geometry}
\usepackage{color}
\usepackage{fancyvrb}
\newcommand{\VerbBar}{|}
\newcommand{\VERB}{\Verb[commandchars=\\\{\}]}
\DefineVerbatimEnvironment{Highlighting}{Verbatim}{commandchars=\\\{\}}
% Add ',fontsize=\small' for more characters per line
\usepackage{framed}
\definecolor{shadecolor}{RGB}{248,248,248}
\newenvironment{Shaded}{\begin{snugshade}}{\end{snugshade}}
\newcommand{\AlertTok}[1]{\textcolor[rgb]{0.94,0.16,0.16}{#1}}
\newcommand{\AnnotationTok}[1]{\textcolor[rgb]{0.56,0.35,0.01}{\textbf{\textit{#1}}}}
\newcommand{\AttributeTok}[1]{\textcolor[rgb]{0.13,0.29,0.53}{#1}}
\newcommand{\BaseNTok}[1]{\textcolor[rgb]{0.00,0.00,0.81}{#1}}
\newcommand{\BuiltInTok}[1]{#1}
\newcommand{\CharTok}[1]{\textcolor[rgb]{0.31,0.60,0.02}{#1}}
\newcommand{\CommentTok}[1]{\textcolor[rgb]{0.56,0.35,0.01}{\textit{#1}}}
\newcommand{\CommentVarTok}[1]{\textcolor[rgb]{0.56,0.35,0.01}{\textbf{\textit{#1}}}}
\newcommand{\ConstantTok}[1]{\textcolor[rgb]{0.56,0.35,0.01}{#1}}
\newcommand{\ControlFlowTok}[1]{\textcolor[rgb]{0.13,0.29,0.53}{\textbf{#1}}}
\newcommand{\DataTypeTok}[1]{\textcolor[rgb]{0.13,0.29,0.53}{#1}}
\newcommand{\DecValTok}[1]{\textcolor[rgb]{0.00,0.00,0.81}{#1}}
\newcommand{\DocumentationTok}[1]{\textcolor[rgb]{0.56,0.35,0.01}{\textbf{\textit{#1}}}}
\newcommand{\ErrorTok}[1]{\textcolor[rgb]{0.64,0.00,0.00}{\textbf{#1}}}
\newcommand{\ExtensionTok}[1]{#1}
\newcommand{\FloatTok}[1]{\textcolor[rgb]{0.00,0.00,0.81}{#1}}
\newcommand{\FunctionTok}[1]{\textcolor[rgb]{0.13,0.29,0.53}{\textbf{#1}}}
\newcommand{\ImportTok}[1]{#1}
\newcommand{\InformationTok}[1]{\textcolor[rgb]{0.56,0.35,0.01}{\textbf{\textit{#1}}}}
\newcommand{\KeywordTok}[1]{\textcolor[rgb]{0.13,0.29,0.53}{\textbf{#1}}}
\newcommand{\NormalTok}[1]{#1}
\newcommand{\OperatorTok}[1]{\textcolor[rgb]{0.81,0.36,0.00}{\textbf{#1}}}
\newcommand{\OtherTok}[1]{\textcolor[rgb]{0.56,0.35,0.01}{#1}}
\newcommand{\PreprocessorTok}[1]{\textcolor[rgb]{0.56,0.35,0.01}{\textit{#1}}}
\newcommand{\RegionMarkerTok}[1]{#1}
\newcommand{\SpecialCharTok}[1]{\textcolor[rgb]{0.81,0.36,0.00}{\textbf{#1}}}
\newcommand{\SpecialStringTok}[1]{\textcolor[rgb]{0.31,0.60,0.02}{#1}}
\newcommand{\StringTok}[1]{\textcolor[rgb]{0.31,0.60,0.02}{#1}}
\newcommand{\VariableTok}[1]{\textcolor[rgb]{0.00,0.00,0.00}{#1}}
\newcommand{\VerbatimStringTok}[1]{\textcolor[rgb]{0.31,0.60,0.02}{#1}}
\newcommand{\WarningTok}[1]{\textcolor[rgb]{0.56,0.35,0.01}{\textbf{\textit{#1}}}}
\usepackage{graphicx}
\makeatletter
\def\maxwidth{\ifdim\Gin@nat@width>\linewidth\linewidth\else\Gin@nat@width\fi}
\def\maxheight{\ifdim\Gin@nat@height>\textheight\textheight\else\Gin@nat@height\fi}
\makeatother
% Scale images if necessary, so that they will not overflow the page
% margins by default, and it is still possible to overwrite the defaults
% using explicit options in \includegraphics[width, height, ...]{}
\setkeys{Gin}{width=\maxwidth,height=\maxheight,keepaspectratio}
% Set default figure placement to htbp
\makeatletter
\def\fps@figure{htbp}
\makeatother
\setlength{\emergencystretch}{3em} % prevent overfull lines
\providecommand{\tightlist}{%
  \setlength{\itemsep}{0pt}\setlength{\parskip}{0pt}}
\setcounter{secnumdepth}{-\maxdimen} % remove section numbering
\ifLuaTeX
  \usepackage{selnolig}  % disable illegal ligatures
\fi
\IfFileExists{bookmark.sty}{\usepackage{bookmark}}{\usepackage{hyperref}}
\IfFileExists{xurl.sty}{\usepackage{xurl}}{} % add URL line breaks if available
\urlstyle{same}
\hypersetup{
  pdftitle={Various\_Simulated\_Data\_Generation\_Setup},
  pdfauthor={RBVNF},
  hidelinks,
  pdfcreator={LaTeX via pandoc}}

\title{Various\_Simulated\_Data\_Generation\_Setup}
\author{RBVNF}
\date{2026-07-30}

\begin{document}
\maketitle

\hypertarget{the-following-is-the-procedure-to-obtain-simulated-data-set.}{%
\section{The following is the procedure to obtain simulated data
set.}\label{the-following-is-the-procedure-to-obtain-simulated-data-set.}}

The Packages to be loaded.

\begin{Shaded}
\begin{Highlighting}[]
\FunctionTok{library}\NormalTok{(RBVNF)}
\end{Highlighting}
\end{Shaded}

\begin{verbatim}
## 
## Attaching package: 'RBVNF'
\end{verbatim}

\begin{verbatim}
## The following object is masked from 'package:base':
## 
##     norm
\end{verbatim}

\begin{Shaded}
\begin{Highlighting}[]
\NormalTok{RBVNF}\SpecialCharTok{::}\FunctionTok{load\_packages}\NormalTok{()}
\end{Highlighting}
\end{Shaded}

\begin{verbatim}
## Loading required package: numDeriv
\end{verbatim}

\begin{verbatim}
## Loading required package: MASS
\end{verbatim}

\begin{verbatim}
## Warning: package 'ggplot2' was built under R version 4.3.3
\end{verbatim}

\begin{verbatim}
## Loading required package: Rcpp
\end{verbatim}

\begin{verbatim}
## Loading required package: RcppZiggurat
\end{verbatim}

\begin{verbatim}
## Loading required package: RcppParallel
\end{verbatim}

\begin{verbatim}
## 
## Attaching package: 'RcppParallel'
\end{verbatim}

\begin{verbatim}
## The following object is masked from 'package:Rcpp':
## 
##     LdFlags
\end{verbatim}

\begin{verbatim}
## 
## Rfast: 2.1.0
\end{verbatim}

\begin{verbatim}
##  ___ __ __ __ __    __ __ __ __ __ _             _               __ __ __ __ __     __ __ __ __ __ __   
## |  __ __ __ __  |  |  __ __ __ __ _/            / \             |  __ __ __ __ /   /__ __ _   _ __ __\  
## | |           | |  | |                         / _ \            | |                        / /          
## | |           | |  | |                        / / \ \           | |                       / /          
## | |           | |  | |                       / /   \ \          | |                      / /          
## | |__ __ __ __| |  | |__ __ __ __           / /     \ \         | |__ __ __ __ _        / /__/\          
## |    __ __ __ __|  |  __ __ __ __|         / /__ _ __\ \        |_ __ __ __ _   |      / ___  /           
## |   \              | |                    / _ _ _ _ _ _ \                     | |      \/  / /       
## | |\ \             | |                   / /           \ \                    | |         / /          
## | | \ \            | |                  / /             \ \                   | |        / /          
## | |  \ \           | |                 / /               \ \                  | |       / /          
## | |   \ \__ __ _   | |                / /                 \ \     _ __ __ __ _| |      / /          
## |_|    \__ __ __\  |_|               /_/                   \_\   /_ __ __ __ ___|      \/             team
\end{verbatim}

\begin{verbatim}
## Loading required package: cowplot
\end{verbatim}

\begin{verbatim}
## 
## Attaching package: 'mvtnorm'
\end{verbatim}

\begin{verbatim}
## The following objects are masked from 'package:Rfast':
## 
##     Crossprod, dmvnorm, dmvt, rmvnorm, rmvt, Tcrossprod
\end{verbatim}

\begin{verbatim}
## Loading required package: Matrix
\end{verbatim}

\begin{verbatim}
## Loaded glmnet 4.1-8
\end{verbatim}

\begin{Shaded}
\begin{Highlighting}[]
\NormalTok{RBVNF}\SpecialCharTok{::}\FunctionTok{load\_additional\_packages}\NormalTok{()}
\end{Highlighting}
\end{Shaded}

\begin{verbatim}
## 
## Attaching package: 'plotly'
\end{verbatim}

\begin{verbatim}
## The following object is masked from 'package:ggplot2':
## 
##     last_plot
\end{verbatim}

\begin{verbatim}
## The following object is masked from 'package:MASS':
## 
##     select
\end{verbatim}

\begin{verbatim}
## The following object is masked from 'package:stats':
## 
##     filter
\end{verbatim}

\begin{verbatim}
## The following object is masked from 'package:graphics':
## 
##     layout
\end{verbatim}

\begin{verbatim}
## Loading required package: dplyr
\end{verbatim}

\begin{verbatim}
## 
## Attaching package: 'dplyr'
\end{verbatim}

\begin{verbatim}
## The following object is masked from 'package:gridExtra':
## 
##     combine
\end{verbatim}

\begin{verbatim}
## The following object is masked from 'package:Rfast':
## 
##     nth
\end{verbatim}

\begin{verbatim}
## The following object is masked from 'package:MASS':
## 
##     select
\end{verbatim}

\begin{verbatim}
## The following objects are masked from 'package:stats':
## 
##     filter, lag
\end{verbatim}

\begin{verbatim}
## The following objects are masked from 'package:base':
## 
##     intersect, setdiff, setequal, union
\end{verbatim}

\begin{verbatim}
## -- Attaching core tidyverse packages ------------------------ tidyverse 2.0.0 --
## v forcats   1.0.0     v stringr   1.5.0
## v lubridate 1.9.3     v tibble    3.2.1
## v purrr     1.0.2     v tidyr     1.3.0
## v readr     2.1.4     
## -- Conflicts ------------------------------------------ tidyverse_conflicts() --
## x dplyr::combine()    masks gridExtra::combine()
## x tidyr::expand()     masks Matrix::expand()
## x dplyr::filter()     masks plotly::filter(), stats::filter()
## x purrr::is_integer() masks Rfast::is_integer()
## x dplyr::lag()        masks stats::lag()
## x dplyr::nth()        masks Rfast::nth()
## x tidyr::pack()       masks Matrix::pack()
## x dplyr::select()     masks plotly::select(), MASS::select()
## x lubridate::stamp()  masks cowplot::stamp()
## x purrr::transpose()  masks Rfast::transpose()
## x tidyr::unpack()     masks Matrix::unpack()
## i Use the conflicted package (<http://conflicted.r-lib.org/>) to force all conflicts to become errors
\end{verbatim}

\hypertarget{simulation-setup-i-simulated-data-when-beta-is-not-sparse.}{%
\section{\texorpdfstring{Simulation Setup I: Simulated data when
\(\beta\) is not
Sparse.}{Simulation Setup I: Simulated data when \textbackslash beta is not Sparse.}}\label{simulation-setup-i-simulated-data-when-beta-is-not-sparse.}}

In the following simulation. Each coordinate of \(\beta^{\tiny{True}}\)
is generated from Uniform distribution between
\([-\text{beta_factor}, \text{beta_factor} ]\). Default value for
beta\_factor=10. Therefore, as a default, the regression coefficients
are generated from \(\text{Uniform}[-10, 10]\). In the demonstration
next, we have used \(p=5\), \(n=500\) and \(d=3\) where \(p\) is number
of covariates, \(n\) is sample size and \(d\) is such that the
\({\bf Y} \in \mathbb{S}^{d-1}\). Therefore, in the following
simulation, \({\bf X} \in \mathbb{R}^{n\times p}\),
\(\beta^{\tiny{True}}\in {\mathbb{R}}^{p\times d}\), and \({\bf Y}\) is
a \(n\times d\) dimensional matrix with unit norm for each of its rows.
In the following simulated data generation code `beta' is same as
\(\beta^{\tiny{True}}\).

\begin{Shaded}
\begin{Highlighting}[]
\NormalTok{n}\OtherTok{=}\DecValTok{500}  \CommentTok{\# Number of the samples}
\NormalTok{p}\OtherTok{=}\DecValTok{5}    \CommentTok{\# Number of the regression covariates}
\NormalTok{d}\OtherTok{=}\DecValTok{3}     \CommentTok{\# Number of direcions in the direcional data}


\DocumentationTok{\#\#\#\# bbeta is a matrix of dimension p\textbackslash{}times d}
\CommentTok{\#bbeta=matrix( rnorm(p*d), nrow=p, ncol=d)}


\NormalTok{data\_lst }\OtherTok{=} \FunctionTok{Data\_generator\_vnf\_reg}\NormalTok{(}\AttributeTok{n=}\NormalTok{n,}
                                  \AttributeTok{p=}\NormalTok{p, }
                                  \AttributeTok{d=}\NormalTok{d, }
                                  \AttributeTok{concentration\_factor =} \DecValTok{1}\NormalTok{, }
                                  \AttributeTok{beta\_factor =} \DecValTok{10}\NormalTok{)}
\end{Highlighting}
\end{Shaded}

\hypertarget{specification-of-the-simulated-data}{%
\subsection{Specification of the simulated
data}\label{specification-of-the-simulated-data}}

\begin{verbatim}
## [1] "Dimension of Y:"
\end{verbatim}

\begin{verbatim}
## [1] 500   3
\end{verbatim}

\begin{verbatim}
## [1] "Dimension of X:"
\end{verbatim}

\begin{verbatim}
## [1] 500   5
\end{verbatim}

\begin{verbatim}
## [1] "Dimension of beta:"
\end{verbatim}

\begin{verbatim}
## [1] 5 3
\end{verbatim}

\begin{verbatim}
## [1] "A Few first rows of Beta"
\end{verbatim}

\begin{verbatim}
##           [,1]      [,2]       [,3]
## [1,] -1.183830 -1.459454  0.7792993
## [2,]  5.850084  6.119266 -1.8867967
## [3,] -7.263654 -4.097099 -5.5875414
## [4,] -6.450984 -2.991619  0.6078511
## [5,] -2.584431  2.785591  2.6198870
\end{verbatim}

\hypertarget{simulation-setup-ii-simulated-data-when-beta-is-sparse.}{%
\section{\texorpdfstring{Simulation Setup II: Simulated data when
\(\beta\) is
Sparse.}{Simulation Setup II: Simulated data when \textbackslash beta is Sparse.}}\label{simulation-setup-ii-simulated-data-when-beta-is-sparse.}}

In the following simulation, approximately 2\% (for demonstration we set
2\%, it can be changed according to the need) of the regression
coefficients are nonzero. If \(\beta^{\tiny{True}}_{i,j}\) is a nonzero
regression coefficient, then its magnitude is between 10 and 20,
i.e.~\(\vert \beta^{\tiny{True}}_{i, j} \vert \in [10,20]\) when
\(\beta^{\tiny{True}}_{i, j}\) is non zero. In the demonstration next,
we have used \(p=200\), \(n=100\) and \(d=3\) where \(p\) is number of
covariates, \(n\) is sample size and \(d\) is such that the
\({\bf Y} \in \mathbb{S}^{d-1}\). Therefore, in the following
simulation, \({\bf X} \in \mathbb{R}^{n\times p}\),
\(\beta^{\tiny{True}}\in {\mathbb{R}}^{p\times d}\), and \({\bf Y}\) is
a \(n\times d\) dimensional matrix with unit norm for each of its rows.
In the following simulated data generation code `beta' is same as
\(\beta^{\tiny{True}}\).

\begin{Shaded}
\begin{Highlighting}[]
\NormalTok{n}\OtherTok{=}\DecValTok{100}  \CommentTok{\# Number of the samples}
\NormalTok{p}\OtherTok{=}\DecValTok{200}    \CommentTok{\# Number of the regression covariates}
\NormalTok{d}\OtherTok{=}\DecValTok{3}  \CommentTok{\# Number of directions in the directional data}

\NormalTok{Num\_of\_nonzero\_beta}\OtherTok{=} \FunctionTok{round}\NormalTok{(p}\SpecialCharTok{*}\NormalTok{.}\DecValTok{02}\NormalTok{)}
\NormalTok{Min\_Non\_Zero\_beta }\OtherTok{=} \DecValTok{10}\NormalTok{ ; Max\_Non\_Zero\_beta }\OtherTok{=} \DecValTok{20} \CommentTok{\#. it means abs(\textbackslash{}beta\_\{i,j\}) is between [10,20] if it is non{-}zero.}

\DocumentationTok{\#\#\#\# bbeta is a matrix of dimension p\textbackslash{}times d}
\CommentTok{\#bbeta=matrix( rnorm(p*d), nrow=p, ncol=d)}
\NormalTok{sigma\_square}\OtherTok{=}\DecValTok{1}
\NormalTok{tau\_square}\OtherTok{=}\DecValTok{1000}

\NormalTok{data\_lst}\OtherTok{\textless{}{-}}\FunctionTok{Data\_generator\_vnf\_reg\_sparse}\NormalTok{(}\AttributeTok{n=}\NormalTok{n,}
                                        \AttributeTok{p=}\NormalTok{p,}
                                        \AttributeTok{d=}\NormalTok{d,}
                                        \AttributeTok{SetUp =} \DecValTok{3}\NormalTok{,}
                                        \AttributeTok{NumOfNonZeroBeta=}\FunctionTok{c}\NormalTok{(Num\_of\_nonzero\_beta,}
\NormalTok{                                                           Min\_Non\_Zero\_beta,}
\NormalTok{                                                           Max\_Non\_Zero\_beta}
\NormalTok{                                        )}
\NormalTok{)}
\end{Highlighting}
\end{Shaded}

\begin{verbatim}
## [1] " Currently Using Default SetUp=3 "
\end{verbatim}

\hypertarget{specification-of-the-simulated-data-1}{%
\subsection{Specification of the simulated
data}\label{specification-of-the-simulated-data-1}}

\begin{verbatim}
## [1] "Dimension of Y:"
\end{verbatim}

\begin{verbatim}
## [1] 100   3
\end{verbatim}

\begin{verbatim}
## [1] "Dimension of X:"
\end{verbatim}

\begin{verbatim}
## [1] 100 200
\end{verbatim}

\begin{verbatim}
## [1] "Dimension of beta:"
\end{verbatim}

\begin{verbatim}
## [1] 200   3
\end{verbatim}

\begin{verbatim}
## [1] "A Few first rows of Beta"
\end{verbatim}

\begin{verbatim}
##           [,1]     [,2]      [,3]
## [1,] -13.61817 18.06304 -18.48712
## [2,] -15.68700 15.53482 -13.09657
## [3,]  11.70414 12.96203  17.73042
## [4,]  19.05281 12.19159  12.36206
## [5,]   0.00000  0.00000   0.00000
## [6,]   0.00000  0.00000   0.00000
\end{verbatim}

\end{document}
